% soundness.tex - Soundness Theorem and Proof

\section{Soundness Theorem}
\label{sec:soundness}

This section explains how the soundness of the capability model is established. %

\begin{theorem}[Soundness]
\label{thm:soundness}
For any capability declaration $\cpdecl$, term $t$, and type $T$:
\[
\typ{\cpdecl}{\cdot}{t}{T}{\effempty}
\quad\Longrightarrow\quad
\exists v.\; \eval{t}{\cdot}{\cpempty}{v}
\]
\end{theorem}

In words: a closed term typed with empty capabilities evaluates successfully
with no capabilities provided. %

\subsection{Why This Is Enough}

This theorem establishes the central static guarantee of the calculus: the
capability annotation in the type system is sound. %
If the type system says a term needs no capabilities, then evaluation agrees. %

The more general statement, used in the proof, is that a term typed with
capability set $\varphi$ evaluates successfully under any capability environment
that supplies $\varphi$. %

\subsection{Proof Strategy}

The proof uses logical relations, a standard technique for semantic type
soundness arguments \cite{pierce-atapl}. %
Rather than reasoning directly about all evaluations, we define semantic
predicates describing when:
\begin{itemize}[nosep]
\item a value satisfies a type,
\item a value environment satisfies a typing environment, and
\item a capability environment satisfies a capability set. %
\end{itemize}

The central proof obligation is then the familiar one: well-typed terms in
satisfying environments evaluate to satisfying values. %

\subsection{The Fundamental Lemma}

\begin{definition}[Logical Relations]
\label{def:logical-relations}
For a capability declaration $\cpdecl$, the proof defines three mutually
supporting semantic predicates. %

\[
\begin{array}{rcl}
\haltsval{\cpdecl}{v}{T}
&\text{means}&
v \text{ is a well-behaved semantic value of type } T, \\[0.25em]
\haltsenv{\cpdecl}{\tpenv}{\vlenv}
&\text{means}&
\vlenv \text{ maps every } x{:}T \in \tpenv
\text{ to some } v \text{ with } \haltsval{\cpdecl}{v}{T}, \\[0.25em]
\haltscpenv{\cpdecl}{\varphi}{\cpenv}
&\text{means}&
\cpenv \text{ supplies every } c \in \varphi
\text{ with a value satisfying } \cpdecl(c).
\end{array}
\]

The value relation is defined by recursion on types. %
Natural-number values satisfy $\tnat$. %
A closure satisfies $\tarr{T_1}{T_2}{\varphi}{\varphi}$ when, for every argument
satisfying $T_1$ and every caller capability environment satisfying $\varphi$,
evaluating the closure body under the extended value environment and merged
capability environment produces a result satisfying $T_2$. %
\end{definition}

\begin{lemma}[Fundamental Lemma]
\label{lem:halts-env-closed}
For any $\cpdecl$, $\tpenv$, $\vlenv$, $\cpenv$, $t$, $T$, and $\varphi$:
\[
\begin{array}{l}
\haltsenv{\cpdecl}{\tpenv}{\vlenv} \;\land \\
\haltscpenv{\cpdecl}{\varphi}{\cpenv} \;\land \\
\typ{\cpdecl}{\tpenv}{t}{T}{\varphi}
\end{array}
\quad\Longrightarrow\quad
\exists v.\; \eval{t}{\vlenv}{\cpenv}{v} \land \haltsval{\cpdecl}{v}{T}
\]
\end{lemma}

This lemma is the real engine of the proof. %
In words, if the variables and capabilities required by a well-typed term are
provided by semantically valid environments, then evaluation succeeds, uses only
the capabilities named by the term's capability annotation, and returns a value
of the stated type. %

For security, the lemma says that the capability annotation is a sound upper
bound on runtime authority: code typed with capability set $\varphi$ cannot
obtain untracked authority outside $\varphi$. %
The theorem follows
immediately by instantiating it with empty environments and empty
capabilities. %

\subsection{Mechanization}

The development is mechanized in Coq. %
The mechanization covers the definitions, auxiliary lemmas on capability
environments, the fundamental lemma, and the final soundness theorem. %

For this paper, the important point is not merely that the theorem has been
checked by a proof assistant. %
The machine-checked proof confirms the mathematical soundness of the capability
model. %
